\documentclass[a4paper, titlepage]{book}

\usepackage{amsmath}
\numberwithin{equation}{section}
\usepackage{amssymb}
\usepackage{csquotes}
\MakeOuterQuote{"}
\usepackage[useregional]{datetime2}
\DTMsetstyle{iso}
\usepackage{hyperref}
\hypersetup{
    colorlinks = true,
    allcolors = blue
}
\usepackage{indentfirst}
\usepackage[skip=5pt, indent=30pt]{parskip}
\usepackage{pgf}
\usepackage{physics2}
\usepackage{siunitx}
\sisetup{
    input-digits = 0123456789\pi,
    inter-unit-product = \cdot,
}


\title{
    Eric Zhang's Physics Notes
}
\author{Eric Zhang}
\date{\today}


\begin{document}
\frontmatter
\maketitle
\tableofcontents


\chapter{Preface}

This document was created as a collection and expansion of my physics notes
from AP Physics C: Electricity and Magnetism, AP Physics C: Mechanics,
the University of Maryland's PHYS273: Intermediate Oscillations and Waves.
Note that the classes are listed in the order that I took them.
However, the notes are arranged with electricity and magnetism \emph{after} mechanics,
since that is the order that most people take them in.

As a result of electricity and magnetism being my first physics class,
the "Electricity and Magnetism" section is relatively
compared to the "Mechanics" and "Waves and Oscillations" section.

These notes are structured in a way to separate physics and mathematical techniques.
Each chapter contains a list of the recommended mathematical topics,
including links to the relevant sections.

\mainmatter
\chapter{Mechanics}

\chapter{Electricity and Magnetism}

Recommended mathematical topics: \ref{section:vectors}

\section{The Electric Force}

\textbf{Electric charge} is denoted by the symbol \(q\) or sometimes \(Q\),
and its SI unit is the coulomb (C).

Charge can be either positive or negative,
with protons carrying a positive charge and electrons carrying a negative charge.

The \textbf{Law of Conservation of Charge} states that charge can not created or destroyed,
only transferred,
and the net charge has to remain constant.
There are two ways to charge and discharge (i.e., move charges):
\begin{itemize}
    \item Conduction -- friction (triboelectricity) or contact
    \item Induction -- contactless
\end{itemize}

Charged objects exert an \textbf{electric force} on other charged objects around them.
The magnitude of the electric force can be given by \textbf{Coulomb's Law}:
\begin{equation}
    F = k\frac{q_1 q_2}{r^2}
\end{equation}
where
\begin{itemize}
    \item \(F\) is the force \emph{both} objects experience as given by Newton's Third Law.
    \item \(k \approx \qty{8.99e-9}{\newton \meter\squared \per \coulomb\squared}\)
          and is called the Coulomb constant.
    \item \(q_1\) and \(q_2\) are the two charges of the two objects.
          Note that both values are signed.
    \item \(r\) is the distance between the two particles.
\end{itemize}
Objects with charges of the same sign will repel each other,
while objects with charges of opposite signs will attract each other.
Thus, Coulomb's Law will give a \emph{negative} value for attraction
and a \emph{positive} value for repulsion.

\section{Maxwell's Equations}


\chapter{Waves and Oscillations} \label{ch:waves-and-oscillations}

Recommended mathematical topics: \ref{section:complex-and-imaginary-numbers}

\section{}

\chapter{Special and General Relativity}

\section{Special Relativity}

\section{General Relativity}


\chapter{Quantum Mechanics}

\section{Schrödinger's Equation}


\appendix
\chapter{Mathematics and Notation}

\section{Vectors} \label{section:vectors}

\section{Complex and Imaginary Numbers} \label{section:complex-and-imaginary-numbers}

\(i = \sqrt{-1}\) is the \textbf{imaginary unit}.

\textbf{Imaginary numbers} are the set of all numbers that can be written in the form \(bi\),
where \(b \in \mathbb{R}\).

\textbf{complex numbers}, denoted by \(\mathbb{C}\),
are the set of all numbers that can be written in the form \(a + bi\),
where \(a,~b \in \mathbb{R}\).

\(a\) is known as the \textbf{real part},
and \(b\) is known as the \textbf{imaginary part}.
Note that this means the imaginary part is \emph{real}.
Given the complex number \(z = a + bi\),
\(\mathrm{Re}[z] = \Re[z] = a\), and \(\mathrm{Im}[z] = \Im[z] = b\).

The \textbf{complex conjugate} of \(z = a + bi\) is given by \(z^* = a - bi\).

\subsection{Complex Plane}

Complex numbers can be represented as 2D vectors on a plane,
with the x-axis being the real part (i.e., \(a\)),
and the y-axis being the imaginary part (i.e., \(b\)).

Given a complex number \(z\),
its magnitude \(\left\|z\right\|\) is given by Pythagorean's Theorem:
\begin{equation}
    \left\|z\right\| = \sqrt{\Re[z]^2 + \Im[z]^2} = \sqrt{z z^*}
\end{equation}

Its angle \(\phi\) is derived through trigonometry
(be careful about the general limitations of \(\arctan\) on calculators):
\begin{equation}
    \phi = \arctan\left(\frac{\Im[z]}{\Re[z]}\right)
\end{equation}

\subsection{Euler's Formula}

Euler's formula:
\begin{equation}
    e^{i \theta} = \cos\theta + i \sin\theta
\end{equation}

\begin{equation}
    \therefore e^{a + bi} = e^a e^{bi} = e^a \left(\cos b + i \sin b\right)
\end{equation}

Oftentimes, working with the complex exponential is easier than the trigonometric functions.

\subsection{Complex-Valued Function}

A \textbf{complex-valued function} is a function that returns a complex number.
It is denoted with a overhead tilde, e.g., \(\tilde{f}\left(x\right)\).

\backmatter
\chapter{References}



\end{document}